% %==============================================================================
% %  CHAPTER 3 — MODEL STRUCTURE AND DESIGN
% %------------------------------------------------------------------------------
% %  Style: matches introduction.tex conventions (\citep/\citet authoryear; keys
% %  resolve against introduction_refs.bib + model_refs.bib). Calibration of the
% %  behavioural parameters is deliberately deferred to the Data & Calibration
% %  chapter; this chapter fixes the model's structure and cites the design
% %  source for every mechanism (Simudyne ODD first; deviations flagged inline).
% %==============================================================================
% The model is a discrete-time agent-based simulation of a centrally cleared
% futures market, built in two coupled tiers. The first tier is a
% \emph{market-microstructure layer}: a limit order book (LOB) on which
% heterogeneous traders exchange a single asset --- the E-mini S\&P~500 (ES)
% front-month future --- around an exogenous fundamental-value process derived
% from real ES data. The second tier is a \emph{central-clearing layer} placed
% on top of the first: a central counterparty (CCP), banking and non-banking
% clearing members, and a population of cleared end-clients, connected by the
% star topology of a real clearing system and governed by
% regulation-grounded margin, default-fund, and default-management mechanics.
% The architecture follows the Simudyne CCP Risk Model
% \citep{simudyne_odd,deloitte_ccp}, which this thesis extends in one central
% direction: the ODD's clearing members trade only for their own account,
% whereas here most market participants do not face the CCP directly but clear
% \emph{through} a member --- the client-clearing tier whose stress behaviour
% is the object of study. Throughout, the design rule is that every mechanism
% and constant is taken from a reference model, from regulation, or measured
% from data; where the implementation deviates from the primary reference, the
% deviation is stated and motivated.

% The two tiers are deliberately asymmetric in how their parameters are set.
% The market layer's behavioural parameters are \emph{calibrated} so that the
% simulated price reproduces the empirical stylised facts of ES returns
% \citep{cont_2001}; the calibration methodology and results are the subject
% of Chapters~4 and~5. The clearing layer contains no free behavioural
% parameters at all: margin rates, default-fund sizing, capital floors, and
% waterfall structure are pinned to their regulatory or reference values. The
% intention is that the clearing tier acts in a market environment that has
% been \emph{earned} from data rather than assumed, while the clearing
% mechanics themselves remain transparent and auditable.

% Time advances in one-minute steps over 6.5-hour regular trading sessions,
% the ODD-native cadence \citep{simudyne_odd}. Two market regimes are used
% throughout: a \emph{calm} regime (calendar year 2019) and a \emph{stressed}
% regime (the COVID-19 crash window, late February to early April 2020), each
% with its own fundamental path and its own calibrated parameter set.

% %==============================================================================
% \section{The market simulation}
% \label{sec:model-market}

% The trading venue is a single limit order book with price--time priority and
% a tick size of \$0.25 (the CME ES specification). Orders accumulate during
% each one-minute step and are crossed in a call auction at the step's end;
% executions print at the resting limit price. The mid-price
% $p_t = (b_t + a_t)/2$ is the average of the best bid and ask, falling back
% to the last traded price when one side of the book is empty. Resting orders
% leave the book only by execution or explicit cancellation: the fundamental
% and momentum traders manage a single standing order each by
% \emph{replace-on-new}, while zero-intelligence orders are cancelled
% stochastically at a per-order rate (Section~\ref{sec:model-agents}). This
% deviates from the ODD's blanket ten-step order lifetime, which was removed
% because it was redundant with --- and masked --- the explicit
% \citet{cont_stoikov_talreja_2008} cancellation rate that is the standard
% order-lifetime mechanism in this model class; the deviation is flagged
% against \citet{simudyne_odd}.

% The runtime population places $100$ agents on the book: $30$ fundamental
% traders, $10$ banking clearing members trading their own account with
% identical fundamental-trader logic, $20$ momentum traders, and $40$
% zero-intelligence traders. Off the book sit five non-banking clearing
% members and the CCP. Ninety of the LOB traders are \emph{cleared clients}
% who access the CCP only through a clearing member
% (Section~\ref{sec:model-margin}).

% %------------------------------------------------------------------------------
% \subsection{Trading agents}
% \label{sec:model-agents}

% Three behavioural types populate the book, the canonical
% fundamentalist--chartist--noise decomposition of the Chiarella family
% \citep{chiarella_1992,majewski_2018} in the LOB form of the Simudyne CCP
% model \citep{simudyne_odd}. All arrival decisions are per-step Bernoulli
% draws at the one-minute cadence, and order sizes are uniform,
% $q \sim \mathcal{U}\{1,\dots,10\}$ \citep{simudyne_odd}.

% \paragraph{Fundamental traders (FT).}
% Each fundamental trader $i$ holds a fixed idiosyncratic belief draw
% $z_i \sim \mathcal{N}(0,1)$ and forms, at every step, a reservation price
% around the fundamental value $V_t$:
% \begin{equation}
% \label{eq:ft-reservation}
% R_{i,t} \;=\; V_t + z_i \,\sigma^{f}_{t},
% \qquad
% \sigma^{f}_{t} \;=\; c_f\,\sigma_t\,v_0 ,
% \end{equation}
% where $\sigma_t$ is the prevailing local volatility of the fundamental
% (Section~\ref{sec:model-fv}), $v_0$ the regime's initial price level, and
% $c_f$ a calibrated scale governing the dispersion of beliefs --- the
% volatility-tracking belief width of the Deloitte specification
% \citep{deloitte_ccp}, which widens the reservation cloud in turbulent
% periods. The trader submits one limit order \emph{at} its reservation price
% on the side implied by its own mispricing,
% $\mathrm{sign}(R_{i,t} - p_{t-1})$, replacing any order it has standing.
% Because $z_i$ is persistent rather than redrawn, optimistic traders
% accumulate long inventories and pessimistic ones short inventories over
% time; this heterogeneous-beliefs specification
% \citep{chiarella_1992,majewski_2018} is load-bearing for the clearing tier,
% since margin calls and default management act on exactly these concentrated
% positions.

% \paragraph{Momentum traders (MT).}
% Each momentum trader tracks an exponentially weighted moving average of
% one-minute mid log-returns $r_t = \ln p_t - \ln p_{t-1}$,
% \begin{equation}
% \label{eq:mt-ewma}
% M_t \;=\; (1-\lambda)\,M_{t-1} + \lambda\, r_{t-1},
% \end{equation}
% and, whenever $|M_t|$ clears a small activation floor, posts a single
% passive limit order on the trend side at a stochastic depth of
% $k \sim \mathrm{Geom}(p_{\mathrm{zi}})$ ticks from the mid, managed
% replace-on-new. The trend horizon $\lambda$ is fixed by design rather than
% estimated, following \citet{majewski_2018}, who likewise impose the
% chartist timescale externally. Momentum traders are limit-only: an earlier
% market-order branch was removed because trend-chasing marketable flow
% corrupted the near-zero autocorrelation of simulated returns.

% \paragraph{Zero-intelligence traders (ZI).}
% The liquidity floor of the book is supplied by zero-intelligence traders in
% the tradition of \citet{farmer_2005} and \citet{cont_stoikov_talreja_2008}.
% In each step a ZI trader independently (i) cancels each of its resting
% orders with probability $\delta_{\mathrm{zi}}$, (ii) submits a limit order
% with probability $\alpha_{\mathrm{zi}}$ at a uniformly random side and a
% geometric depth $k \sim \mathrm{Geom}(p_{\mathrm{zi}})$ from the mid, and
% (iii) submits a market order with probability $\mu_{\mathrm{zi}}$, with no
% dependence on any signal. The placement parameter $p_{\mathrm{zi}}$ is not a
% free parameter: it is fitted by maximum likelihood to the depth distribution
% of new order arrivals in ES market-by-price (MBP-10) data, so the simulated
% book's near-mid thickness is anchored to the measured one. The arrival and
% cancellation rates are calibrated; the market-order rate is pinned at the
% \citet{cont_stoikov_talreja_2008} baseline.

% Two agent types that appear in related designs are deliberately absent. An
% always-quoting market maker \citep{gao_2022_hfabm} and a volatility trader
% \citep{gao_2023} were both implemented and removed after an ablation study
% showed the former damps volatility and clustering in both regimes while the
% latter helps only under stress; the ablation is reported with the
% calibration results.

% %------------------------------------------------------------------------------
% \subsection{The fundamental value process}
% \label{sec:model-fv}

% The fundamental $V_t$ is exogenous to the simulation: a path generated
% offline and read in as a time series, following the data-driven fundamental
% signal of the ODD \citep{simudyne_odd} and the data-derived fundamental of
% the XGB-Chiarella line \citep{gao_2022_xgb}. Two interchangeable generators
% are used, with distinct roles.

% \paragraph{Primary: the observed efficient price.}
% For the headline experiments the fundamental is extracted from the real ES
% one-minute mid by a local-level state-space model,
% \begin{equation}
% \label{eq:kalman}
% \ln V_{t+1} = \ln V_t + \mu + \sigma_v \varepsilon_t,
% \qquad
% \ln p^{\mathrm{obs}}_t = \ln V_t + \eta_t,
% \qquad
% \eta_t \sim \mathcal{N}(0, \sigma_\eta^2),
% \end{equation}
% estimated per session by Kalman maximum likelihood and smoothed by the RTS
% recursion, in the spirit of \citet{gao_2022_xgb} and of
% \citet{majewski_2018}, who treat the fundamental as a latent state filtered
% from price. One property of this construction is reported plainly: at the
% one-minute frequency the ES \emph{mid} contains almost no i.i.d.\
% microstructure noise (the noise that model \eqref{eq:kalman} is built to
% remove lives in trade prices, not in quote midpoints), so the estimated
% $\sigma_\eta$ is of order $10^{-5}$ and the smoother is near-identity. The
% fundamental is therefore, in effect, the observed efficient mid --- the
% Kalman step \emph{verifies} that mid-price noise is negligible rather than
% delivering material smoothing --- and the simulated market inherits the real
% return tails through $V_t$, which the agent layer must transmit and reshape
% rather than manufacture from nothing. Two features of the construction
% matter for the clearing tier. First, the per-session series are concatenated
% at their true levels, so \emph{overnight gaps are retained}: the COVID
% episode carries its actual peak-to-trough drawdown of roughly $-33\%$, most
% of which occurred between sessions. The simulator opens each session by
% repricing the entire book at the gapped fundamental, a clean between-session
% jump that stresses the cleared positions without contaminating intraday
% returns. Second, a per-minute local volatility
% $\sigma_t = \sqrt{\mathrm{EWMA}[r^2_t]}$ accompanies the path and drives the
% fundamental traders' belief width in \eqref{eq:ft-reservation}.

% \paragraph{Scenario generator: a stochastic-volatility jump diffusion.}
% For Monte-Carlo experiments that require many independent stress episodes, a
% synthetic alternative generates $V_t$ as a geometric diffusion whose
% volatility mean-reverts as an Ornstein--Uhlenbeck process
% \citep{stein_stein_1991,heston_1993}, overlaid with compound-Poisson jumps
% \citep{merton_1976} and with empirical overnight gaps:
% \begin{align}
% \sigma_{t+1} &= \theta + e^{-\alpha}(\sigma_t - \theta) + s\,\varepsilon^\sigma_t,
% \label{eq:svmjd-vol}\\
% \ln V_{t+1} - \ln V_t &=
% \begin{cases}
% \mu + \sigma_t \varepsilon_t + J_t, & J_t \sim \mathcal{N}(0,\delta^2)
% \text{ w.p. } \lambda/390, & \text{(intraday steps)}\\[2pt]
% G_t \sim \widehat{F}_{\mathrm{on}}, & & \text{(session boundaries)},
% \end{cases}
% \label{eq:svmjd-ret}
% \end{align}
% where $\widehat{F}_{\mathrm{on}}$ is the empirical distribution of the
% regime's observed overnight returns, resampled nonparametrically. Every
% parameter of \eqref{eq:svmjd-vol}--\eqref{eq:svmjd-ret} is measured from the
% ES data with standard realized-measures estimators rather than assumed: the
% diffusion variance from jump-robust daily bipower variation
% \citep{bns_2004}, the jump intensity from threshold exceedance counts
% \citep{mancini_2009} --- approximately $1.9$ per day in the calm regime and
% $1.0$ per day under stress, superseding the FTSE-based intensity of the ODD,
% a flagged deviation --- and the jump size from the realized-variance minus
% bipower-variation jump component, whose share of intraday variance
% ($6.8\%$ calm, $3.9\%$ stressed) sits on the $5$--$7\%$ benchmark
% established for the S\&P~500 \citep{huang_tauchen_2005}. The
% Ornstein--Uhlenbeck parameters are fitted to a jump-robust intraday
% volatility series with noise-robust autocovariance moments, and the
% long-run level $\theta$ is re-anchored so the stochastic-volatility layer
% redistributes, rather than inflates, the measured variance. The overnight
% component is included because gaps carry roughly $40\%$ of ES close-to-close
% variance in both regimes and because the COVID decline itself was
% concentrated overnight while intraday returns netted positive --- the
% overnight/intraday asymmetry documented by \citet{lou_polk_skouras_2019}.
% Stressed episodes default to the empirical window length of $29$ sessions,
% so that a ``stressed scenario'' means an episode of COVID-like length and
% severity rather than an indefinitely compounded crash.

% %==============================================================================
% \section{The margin framework}
% \label{sec:model-margin}

% The clearing tier connects the market to a single CCP through the star
% topology of \citet{simudyne_odd}: every clearing member holds a bilateral
% link to the CCP, which becomes the counterparty to all cleared positions
% upon novation. Two member types are distinguished, as in the ODD. A
% \emph{banking clearing member} (BCM) trades its own account --- with
% fundamental-trader logic, so the market-layer calibration is unaffected by
% the cast --- and additionally guarantees a book of clients. A
% \emph{non-banking clearing member} (NBCM) is a pure intermediary with no
% proprietary position; its entire balance-sheet capacity is client clearing.
% The thesis extension is the client tier itself: ninety end-clients
% (fundamental, momentum, and zero-intelligence types) clear through the ten
% client-carrying members in skewed books of five to fifteen clients, with the
% non-banking members carrying the largest books. Loss-absorbing capital is
% initialised at regulatory scale --- bank-level capital for BCMs, the
% non-bank futures-commission-merchant range for NBCMs, and entity-type-scaled
% cash for clients --- taken from CFTC member financial data and CCP
% disclosures.

% \paragraph{Variation margin.}
% Every 60 minutes the margin cycle marks each cleared book to the prevailing
% mid and settles the move in cash. Settlement is computed against
% \emph{average filled prices}, the ODD's margin-call profit-and-loss
% convention \citep{simudyne_odd}: writing $N_\tau$ for the end-of-cycle
% position, $m_\tau$ for the mark, and $\{(q_k, p_k)\}$ for the fills executed
% during the cycle,
% \begin{equation}
% \label{eq:vm}
% \mathrm{VM}_\tau \;=\;
% \Bigl[\, N_\tau\,(m_\tau - m_{\tau-1})
% \;+\; \textstyle\sum_k q_k\,(m_{\tau-1} - p_k) \Bigr]\cdot \ell \cdot \chi ,
% \end{equation}
% where $\ell$ is the per-regime contract lot that maps model volume to
% empirical ES volume and $\chi = \$50$ per index point is the CME multiplier.
% The second term in \eqref{eq:vm} settles each trade from its execution
% price, so that the slippage of forced liquidations
% (Section~\ref{sec:model-default}) is borne in cash by the liquidating party
% rather than lost between marks. Cleared participants use futures-style
% accounting throughout: a fill moves only the position, and cash moves only
% through the margin cycle.

% \paragraph{Initial margin.}
% Initial margin is charged on gross exposure as a value-at-risk scan --- for
% a single linear future, the SPAN scanning charge coincides with parametric
% VaR --- per EMIR's requirements for exchange-traded derivatives
% \citep{emir_rts}:
% \begin{equation}
% \label{eq:im}
% f^{\mathrm{IM}}_\tau \;=\;
% \max\!\Bigl( f_{\mathrm{floor}},\; z_{99}\,\hat\sigma^{\mathrm{day}}_\tau
% \sqrt{h} \Bigr),
% \qquad
% \mathrm{IM}_\tau = f^{\mathrm{IM}}_\tau \cdot
% \bigl(\,|N^{\mathrm{own}}_\tau| + \textstyle\sum_{c}|N^{c}_\tau|\,\bigr)
% \cdot \ell\,\chi\, m_\tau ,
% \end{equation}
% with $z_{99}$ the 99\% normal quantile, $h = 2$ days the margin period of
% risk, and an anti-procyclicality floor $f_{\mathrm{floor}}$ anchored to the
% CME's ES margin scale. The driving volatility $\hat\sigma^{\mathrm{day}}_\tau$
% defines the \emph{margin regime}, the experimental lever for the
% procyclicality hypothesis: in the \emph{reactive} baseline it is an
% exponentially weighted estimate of realised intraday simulation returns
% (half-life two sessions), so margin ratchets upward through a crash exactly
% as CME's ES margin was raised repeatedly through March 2020; a \emph{static}
% variant fixes it at the regime constant; and a \emph{flat} variant removes
% the volatility response altogether, charging a constant fraction --- the
% through-the-cycle comparator contemplated by the anti-procyclicality
% options of EMIR and the post-2020 margin debate
% \citep{murphy_2014,esrb_2020}. A maintenance threshold at $95\%$ of initial
% margin flags margin calls when a cycle's variation margin exceeds the
% remaining buffer \citep{simudyne_odd}.

% \paragraph{Capital adequacy.}
% Each member's capital ratio is cash over initial margin,
% \begin{equation}
% \label{eq:capital}
% \kappa_\tau \;=\; \frac{C_\tau}{\mathrm{IM}_\tau} \;\geq\; 8\%,
% \end{equation}
% the futures-commission-merchant net-capital rule, which sets required
% capital against \emph{risk margin} rather than gross notional
% \citep{cftc_reg117}. The distinction matters operationally: an FCM clears
% many multiples of its capital in client notional, so a notional-based floor
% would bind permanently, whereas \eqref{eq:capital} binds only as distress
% approaches. A banking member's proprietary book is additionally capped by a
% trading-book VaR limit, $|N^{\mathrm{own}}|\,\ell\chi m \leq
% \beta C / (z_{99}\hat\sigma^{\mathrm{day}})$ with $\beta = 5\%$ of capital,
% in the spirit of regulatory trading-book risk limits \citep{basel_frtb};
% without it, the fundamental-trader logic accumulates unbounded prop
% positions in a trend.

% \paragraph{Client clearing.}
% Each client carries its own balance sheet. The margin cycle issues every
% client its own variation-margin call, paid from the client's own cash and
% relayed by its member to the CCP; the client's live inventory is novated
% into the member's cleared book each cycle. A client can open positions only
% up to its broker house margin --- a fixed $20\%$ of notional, i.e.\
% five-times leverage, the conservative house multiple above the exchange
% minimum --- and is frozen (barred from new risk) when its cash falls below
% $8\%$ of its gross notional. Cash exhaustion is a client default and hands
% the position to the member (Section~\ref{sec:model-default}). This
% client tier is the thesis's central extension of the reference model;
% its empirical motivation --- clients supply the majority of margin at major
% CCPs and typically depend on a single clearing agent --- is documented by
% \citet{ofr_2026}.

% %==============================================================================
% \section{The default management framework}
% \label{sec:model-default}

% Default management proceeds in three escalating stages --- client default,
% member distress, member default --- and is closed by the CCP's prefunded
% resources. All liquidations use the same execution schedule: a position of
% size $X$ unwound over horizon $H$ follows the closed-form
% \citet{almgren_chriss_2000} trajectory
% \begin{equation}
% \label{eq:ac}
% \frac{x_j}{X} \;=\; \frac{\sinh\bigl(\kappa (H - j)\bigr)}{\sinh(\kappa H)},
% \qquad j = 0,\dots,H,
% \end{equation}
% with urgency $\kappa H = 2$ and $H = 30$ minutes, executed as market orders
% on the simulated book --- so liquidation impact is endogenous: the fire-sale
% walks the very book on which surviving members are marked, which is the
% price-mediated contagion channel of the model
% \citep{simudyne_odd,vytelingum_2025}.

% \paragraph{Client default.}
% A client whose cash is exhausted by variation margin defaults. Its member
% covers the uncollected margin call, assumes the client's position onto its
% own book, and liquidates it by \eqref{eq:ac} --- a banking member through
% its own orders, a non-banking member through orders routed by the CCP on its
% behalf. Because variation margin settles against fill prices
% \eqref{eq:vm}, the close-out loss is realised as the fire-sale walks the
% price: the member's loss is the deficit the liquidation actually produces,
% not a stylised haircut. This client$\to$member loss assumption is the first
% link of the cascade under study.

% \paragraph{Member distress.}
% A member breaching the capital floor \eqref{eq:capital} while still solvent
% acts before defaulting, asymmetrically by type as in the ODD
% \citep{simudyne_odd}: a banking member \emph{deleverages}, selling its
% proprietary book down by \eqref{eq:ac} until the floor is restored; a
% non-banking member, having no own book to sell, \emph{stops out} ---
% suspending new client risk, which freezes its clients --- for as long as the
% breach persists. The stop-out is state-based: it lifts if the ratio
% recovers. These are the liquidity-withdrawal behaviours that couple margin
% pressure back into the market layer.

% \paragraph{Member default and the waterfall.}
% A member whose cash is exhausted defaults. The CCP absorbs the member's
% realised cash deficit --- the shortfall it owes but cannot pay --- through
% the five-level waterfall of the reference model \citep{simudyne_odd}, in
% EMIR's defaulter-pays ordering \citep{emir_rts}:
% \begin{equation}
% \label{eq:waterfall}
% \text{L1 own DF contribution}
% \;\to\; \text{L2 CCP skin-in-the-game}
% \;\to\; \text{L3 pooled DF}
% \;\to\; \text{L4 survivors' cash}
% \;\to\; \text{L5 CCP capital}.
% \end{equation}
% The defaulted member is then removed from the system (ODD step sequence).
% Its surviving clients are \emph{ported}: each is reassigned to a surviving
% client-carrying member with spare risk capacity, accepted only if the
% receiver remains above its capital floor with the transferred book ---
% porting is the regulatory first option on a member default
% \citep{emir_rts}, but it is capacity-checked rather than guaranteed, since
% in practice the modal client has no standing backup agent
% \citep{ofr_2026}; clients that cannot be placed are closed out and
% deactivated. The member's own book, together with any unported client
% positions, passes to the CCP, which disposes of it by \eqref{eq:ac} on the
% open market. The CCP carries this book in a disposal account marked through
% the margin cycle: disposal-period losses re-enter the waterfall
% \eqref{eq:waterfall} of the originating default, so the mutualised loss is
% the close-out loss actually realised, symmetric with the client treatment.
% The open-market disposal is a flagged deviation from the ODD, whose default
% positions are auctioned to the surviving member with the largest opposing
% position; the deviation is deliberate, as the endogenous price impact of
% disposal is precisely the contagion mechanism the thesis measures.

% \paragraph{The default fund.}
% The mutualised layers of \eqref{eq:waterfall} are prefunded. The fund is
% sized daily to the \emph{cover-two} standard --- sufficient for the
% simultaneous default of the two largest members
% \citep{emir_rts} --- by the stress-loss-over-initial-margin method of
% Euronext Clearing \citep{euronext_a9}:
% \begin{equation}
% \label{eq:df}
% \mathrm{DF}_\tau \;=\; (1+b) \sum_{j \in \text{top-2}}
% \max\Bigl( 0,\; \bigl(s_\tau - f^{\mathrm{IM}}_\tau\bigr)\,
% |N_{j,\tau}|\,\ell\,\chi\,m_\tau \Bigr),
% \qquad
% s_\tau = \max\bigl(15\%,\; z_{s}\,\hat\sigma^{\mathrm{day}}_\tau\sqrt{h}\bigr),
% \end{equation}
% the loss each member would suffer beyond posted margin under an
% extreme-but-plausible move $s_\tau$, with buffer $b = 10\%$. The CCP
% prefunds a skin-in-the-game share ahead of the pooled fund, consumed at L2
% before any survivor resources \citep{simudyne_odd}; the EMIR skin-in-the-game
% rule is expressed against CCP regulatory capital rather than fund size, so
% the model's fund-share convention is a stylised assumption and is disclosed
% as such. Member contributions are paid \emph{in cash} into the fund, pro
% rata to exposure, and every post-draw daily recomputation is therefore a
% replenishment cash call on survivors --- the margin-liquidity channel that
% made the March 2020 episode a funding event as much as a solvency one
% \citep{esrb_2020}. Reverse-stress experiments scale the fundamental path
% until \eqref{eq:df} is breached, the scenario-design logic of
% \citet{euronext_a9}.

% Two structural counterfactuals complete the design and operationalise the
% research questions: a \emph{direct-clearing} configuration, in which the
% client tier is removed and all end-users face the CCP directly with
% otherwise identical trading mechanics --- isolating the loss-absorption
% structure of tiering in the spirit of the netting and tiering trade-offs of
% \citet{duffie_zhu_2011} and \citet{galbiati_soramaki_2013} --- and the
% margin-regime switch of Section~\ref{sec:model-margin}. Both are evaluated
% as paired experiments on identical fundamental paths and random seeds, so
% that differences in outcomes are attributable to the structural change
% alone.